\begin{table}[htbp]
\centering
\small
\begin{threeparttable}
\caption{Classifier performance at the binary cut (easier, A or B1, versus harder, B2 or C), test set}
\label{tab:clf-binary-cut}
\begin{tabular}{lrrrrrrrrr}
\hline
Language & n & Binary accuracy & Recall harder & Recall easier & True easier (\%) & Predicted easier (\%) & Harder as easier & Easier as harder & McNemar p \\
\hline
German & 154 & 0.916 & 0.891 & 0.933 & 58.4 & 59.1 & 7 & 6 & 1.000 \\
English & 139 & 0.942 & 0.954 & 0.923 & 37.4 & 37.4 & 4 & 4 & 1.000 \\
Dutch & 120 & 0.775 & 0.660 & 0.849 & 60.8 & 65.0 & 16 & 11 & 0.442 \\
\hline
All & 413 & 0.884 & 0.864 & 0.902 & 52.1 & 53.5 & 27 & 21 & 0.471 \\
\hline
\end{tabular}
\begin{tablenotes}
\footnotesize
\item Document-level predictions of run cefr\_20260817\_231925 on the 413 test texts (window logits averaged per text, as for the corpus pages). The four-level labels are collapsed to the binary outcome of the regression: easier (A or B1) and harder (B2 or C).
\item The two error columns are the off-diagonal cells of the binary confusion matrix; the predicted share of easier texts exceeds the true share when harder as easier exceeds easier as harder. McNemar p: exact two-sided binomial test of these two counts.
\end{tablenotes}
\end{threeparttable}
\end{table}
